%% sbc_crypto.tex — Cryptographic Grounding for SBC (Issue #86, F8)
%% Paper F: "Physically-Constrained Semantic Security Theory for IoT Mesh Routing"
%% Target venue: IEEE Transactions on Dependable and Secure Computing (IEEE TDSC)
%%
%% Purpose: Anchors the Semantic Binding Condition (Definition 6) in a standard
%%          cryptographic hardness assumption (EUF-CMA), resolving the gap W1
%%          identified in the Le Chun evaluation.
%%
%% Design choice: Option A — reposition SBC as a protocol analysis methodology
%%   that invokes a single crypto assumption (Assumption 1), rather than Option B
%%   (define SBC itself as a security game), which would require 2-3 pages of
%%   game-based machinery inappropriate for the TDSC systems/security audience.
%%
%% Required preamble packages:
%%   \usepackage{amsthm, amsmath, amssymb, mathtools}
%%
%% Theorem environments (declare in preamble if not already present):
%%   \theoremstyle{definition}
%%   \newtheorem{assumption}[definition]{Assumption}
%%   \newtheorem{remark}[definition]{Remark}

% -----------------------------------------------------------------------
\subsection{Cryptographic Grounding of SBC}
\label{subsec:sbc-crypto}
% -----------------------------------------------------------------------

Definition~\ref{def:sbc} (SBC) specifies that every message in
$\mathrm{DC}(P)$ must carry an \emph{unforgeable authenticator}, but
does not name the underlying cryptographic hardness assumption.
The proof of Theorem~\ref{thm:sbc-sufficiency} invokes ``PRF security
of HMAC'' informally (Step~1).  This subsection makes that assumption
explicit and repositions SBC as a \emph{protocol analysis methodology}
that reduces routing security to a single standard cryptographic
assumption.

% -----------------------------------------------------------------------
\subsubsection*{Assumption 1: EUF-CMA Security of the MAC Scheme}
% -----------------------------------------------------------------------

\begin{assumption}[EUF-CMA MAC]
\label{asm:euf-cma}
Let $\lambda \in \mathbb{N}$ be the security parameter and let
$\Pi = (\mathcal{K}\text{gen}, \mathrm{MAC}, \mathrm{Vfy})$ be a
message authentication code scheme.  Consider the following
\emph{existential unforgeability under chosen-message attack}
(EUF-CMA) game between a challenger $\mathcal{C}$ and a probabilistic
polynomial-time (PPT) adversary $\mathcal{A}$:

\begin{enumerate}[label=(\arabic*)]
  \item $\mathcal{C}$ samples $k \xleftarrow{\$} \mathcal{K}\text{gen}(1^\lambda)$
        and gives $1^\lambda$ to $\mathcal{A}$.
  \item $\mathcal{A}$ makes arbitrarily many adaptive \emph{signing oracle}
        queries: for each query $m_i$, $\mathcal{C}$ returns
        $\sigma_i \leftarrow \mathrm{MAC}_k(m_i)$.
  \item $\mathcal{A}$ outputs a forgery candidate $(m^*, \sigma^*)$.
  \item $\mathcal{A}$ \emph{wins} if $\mathrm{Vfy}_k(m^*, \sigma^*) = 1$
        and $m^* \notin \{m_1, m_2, \ldots\}$ (i.e., $m^*$ was never
        submitted to the signing oracle).
\end{enumerate}

We define the \emph{EUF-CMA advantage} of $\mathcal{A}$ as
\[
  \mathrm{Adv}^{\mathrm{euf-cma}}_{\Pi}(\mathcal{A}, \lambda)
  \;=\;
  \Pr\!\bigl[\,\mathrm{EUF\text{-}CMA}(\mathcal{A}, \Pi, \lambda) = 1\,\bigr].
\]

\textbf{Assumption 1} states: $\Pi$ is \emph{EUF-CMA secure} if for
every PPT adversary $\mathcal{A}$,
\[
  \mathrm{Adv}^{\mathrm{euf-cma}}_{\Pi}(\mathcal{A}, \lambda)
  \;\leq\;
  \mathrm{negl}(\lambda),
\]
where $\mathrm{negl}(\lambda)$ denotes any function that decreases
faster than the inverse of any polynomial.

\textbf{Instantiation.}  In this paper, $\Pi$ is instantiated as
HMAC-SHA256 with key length $\lambda = 256$ bits.
HMAC-SHA256 is EUF-CMA secure under the assumption that SHA-256 is a
pseudorandom function family~\cite{bellare1996keying}.
\end{assumption}

\begin{remark}[PRF vs.\ EUF-CMA]
\label{rem:prf-euf}
Theorem~\ref{thm:sbc-sufficiency} (Step~1) invokes ``PRF security
of HMAC.''  This is a \emph{sufficient} condition for EUF-CMA security:
any MAC scheme whose tag function is a PRF is automatically EUF-CMA
secure~\cite{katz2020introduction}.  Assumption~\ref{asm:euf-cma}
unifies both formulations under the weaker, standard EUF-CMA condition,
which is the appropriate security notion for message authentication in
protocol analysis.
\end{remark}

% -----------------------------------------------------------------------
\subsubsection*{SBC as Protocol Analysis Methodology}
% -----------------------------------------------------------------------

\begin{remark}[SBC is a protocol analysis tool, not a security definition]
\label{rem:sbc-methodology}
$\mathrm{SBC}(P)$ (Definition~\ref{def:sbc}) should be understood as a
\emph{three-step analysis methodology}:

\begin{enumerate}[label=\textbf{Step~\arabic*.}, leftmargin=2.5em]
  \item \textbf{Identify $\mathrm{DC}(P)$.}  Enumerate the minimal set of
        protocol messages whose forgery or manipulation is sufficient to
        install an arbitrary next-hop entry at any honest node
        (Definition~\ref{def:decision-critical}).

  \item \textbf{Check completeness.}  Verify that every message in
        $\mathrm{DC}(P)$ carries an authenticator computed under $\Pi$
        (conditions~\ref{sbc:auth} and~\ref{sbc:closure} of
        Definition~\ref{def:sbc}).

  \item \textbf{Invoke Assumption~\ref{asm:euf-cma}.}  If the
        completeness check passes \emph{and} $\Pi$ satisfies
        Assumption~\ref{asm:euf-cma}, then $\mathrm{SBC}(P)$ holds
        and the routing security guarantee of
        Theorem~\ref{thm:sbc-sufficiency} follows.
\end{enumerate}

The three SBC conditions in Definition~\ref{def:sbc} (\ref{sbc:auth},
\ref{sbc:fresh}, \ref{sbc:closure}) are \emph{protocol-level
requirements}, not cryptographic security games.  Their force derives
entirely from Assumption~\ref{asm:euf-cma}: without it, the
requirements are vacuous.  With it, the requirements reduce routing
integrity to the EUF-CMA guarantee of HMAC-SHA256.
\end{remark}

% -----------------------------------------------------------------------
\subsubsection*{Revised Theorem 1 Preamble}
% -----------------------------------------------------------------------

The following is the full, assumption-explicit statement of
Theorem~\ref{thm:sbc-sufficiency}.  (The proof in
Theorem~\ref{thm:sbc-sufficiency} is unchanged; this restatement adds an
explicit preamble that names the assumption.)

\begin{theorem}[SBC Sufficiency for Relay Security — Assumption-Explicit Form]
\label{thm:sbc-sufficiency-explicit}
Let $\Pi = (\mathcal{K}\text{gen}, \mathrm{MAC}, \mathrm{Vfy})$
satisfy \textbf{Assumption~\ref{asm:euf-cma}} (EUF-CMA security of
HMAC-SHA256).  Let $P$ be a mesh routing protocol with key $K$
shared among $\mathrm{AuthorizedSet}(P)$, and let $A$ be a Dolev-Yao
attacker with physically-bounded capabilities: duty cycle $\leq \delta$,
minimum airtime per message $\geq \tau$, and maximum injection rate
$\rho_A \leq \delta/\tau$.

If $\mathrm{SBC}(P)$ holds (Definition~\ref{def:sbc}) and
$A \notin \mathrm{AuthorizedSet}(P)$, then for every protocol execution
trace $\mathbf{t}$ and every source node $S$,
\[
  \mathrm{next\_hop}_{S}(\mathbf{t}) \;\neq\; A,
\]
except with probability at most
\[
  \epsilon(\lambda)
  \;=\;
  q(\lambda) \cdot \mathrm{Adv}^{\mathrm{euf-cma}}_\Pi(\mathcal{A}, \lambda)
  \;\leq\;
  q(\lambda) \cdot \mathrm{negl}(\lambda),
\]
where $q(\lambda)$ is a polynomial upper-bounding the total number of
routing-message forgery attempts by $A$ during the execution (bounded
in practice by the physical injection rate $\rho_A$).
\end{theorem}

\begin{remark}[Tightness of the reduction]
\label{rem:reduction-tightness}
The factor $q(\lambda)$ arises from a standard hybrid argument: $A$ can
make at most $q(\lambda)$ forgery attempts during the trace, each
contributing at most $\mathrm{Adv}^{\mathrm{euf-cma}}_\Pi$ to the
overall success probability via a union bound.  For the IoT mesh
setting (LoRa duty-cycle $\delta = 1\%$, message airtime $\tau \geq 0.2$\,s),
$\rho_A \leq 18$\,msg/min, so $q(\lambda)$ grows only linearly in trace
duration.  At $\lambda = 256$ bits, $\mathrm{negl}(\lambda) \approx 2^{-256}$,
making $\epsilon(\lambda)$ negligible even for week-long deployments.
\end{remark}

% -----------------------------------------------------------------------
\subsubsection*{Remarks on Security Model Scope}
% -----------------------------------------------------------------------

\begin{remark}[SBC vs.\ UC Security]
\label{rem:sbc-vs-uc}
$\mathrm{SBC}(P)$ is a \emph{weaker} notion than Universal
Composability (UC) security~\cite{canetti2001universally}, and this
is by design.

UC security requires a \emph{simulation-based} proof: one must exhibit
an ideal-world simulator $\mathcal{S}$ such that no environment
$\mathcal{Z}$ can distinguish real executions of $P$ from ideal
executions.  UC guarantees \emph{cryptographic composability}: a
protocol that is UC-secure for a functionality $\mathcal{F}$ remains
secure when composed with arbitrary other protocols running
concurrently.

$\mathrm{SBC}(P)$ requires only:
\begin{enumerate}[label=(\roman*)]
  \item EUF-CMA security of the MAC scheme (Assumption~\ref{asm:euf-cma}), and
  \item a protocol-level completeness check on $\mathrm{DC}(P)$.
\end{enumerate}

The weaker assumption is \emph{sufficient} for IoT mesh routing for the
following reason: IoT mesh protocols operate in a single administrative
domain with a pre-shared key, a single active session per node, and no
concurrent session interleaving that a UC adversary could exploit.
Specifically:
\begin{itemize}
  \item LoRa half-duplex radio prevents simultaneous transmission and
        reception; concurrent session exploitation is physically
        impossible at a single node.
  \item The duty-cycle constraint ($\leq 1\%$) limits the attacker's
        message injection rate to $\rho_A \leq 18$\,msg/min, ruling out
        the high-volume adaptive chosen-ciphertext queries that motivate
        UC security in the multi-session setting.
  \item Pre-shared key deployment means there is no key-exchange
        sub-protocol to compose with; the only cryptographic operation
        is HMAC verification.
\end{itemize}

Adopting UC security would add substantial proof machinery (simulator
construction, environment distinguishers) without providing any
additional protection against the physically-constrained Dolev-Yao
adversary modelled here.  We therefore follow the standard practice of
applied protocol analysis~\cite{cremers2019comprehensive} and state the
minimal sufficient assumption (Assumption~\ref{asm:euf-cma}) explicitly.
\end{remark}

\begin{remark}[Why not define SBC as a security game? (Option B)]
\label{rem:why-not-game}
An alternative design (Option~B) would define $\mathrm{SBC}(P)$ as a
formal security game: an adversary is given oracle access to the
routing protocol and wins if it installs an unauthorized next-hop entry.
This would make SBC a self-contained security definition comparable to
those in the cryptographic protocol literature.

We deliberately chose \textbf{Option~A} (SBC as methodology invoking
Assumption~\ref{asm:euf-cma}) for two reasons:

\begin{enumerate}[label=(\arabic*)]
  \item \textbf{Audience.}  IEEE TDSC targets the systems-security
        community.  A 2--3-page game-based definition would shift the
        paper's contribution toward a cryptography-theory result,
        diluting the primary contribution: a framework for \emph{protocol
        analysis} applicable across IoT mesh stacks.  Assumption~\ref{asm:euf-cma}
        is a one-paragraph standard reference; it anchors the theory
        without monopolizing the exposition.

  \item \textbf{Scope precision.}  The gap identified by the Le Chun
        evaluation (W1) is that ``SBC is stated as design requirement
        but does not specify the underlying cryptographic assumption.''
        Option~A closes this gap minimally and precisely: it names the
        assumption (EUF-CMA), cites the standard reference, and
        propagates it into Theorem~\ref{thm:sbc-sufficiency} as an
        explicit hypothesis.  Option~B would over-engineer the fix and
        introduce new proof obligations (game reduction, simulator
        construction) that are unnecessary for the theorem as stated.
\end{enumerate}

Authors who wish to extend the framework to a multi-domain or
multi-session setting (where composability matters) may layer a
UC-style proof on top of Assumption~\ref{asm:euf-cma}; this is left
as future work.
\end{remark}
